\documentclass[a4paper, 12pt] {article}

\usepackage[margin=1in]{geometry}
\usepackage[framemethod=TikZ]{mdframed}
\usepackage{graphicx}
\usepackage{xcolor,colortbl}
\usepackage{listings} % Allows code-listings
\usepackage{wrapfig}
\usepackage{courier} % proper monospace font for code
\usepackage{lscape}
\usepackage[utf8]{inputenc}
\usepackage[norsk]{babel}
\usepackage[T1]{fontenc}
\usepackage{rotating}
\usepackage{epstopdf}
\usepackage{hyperref} %linking fra content  table
\usepackage{xstring} %string-manipulation
\usepackage{enumitem} %used for enumerate-manipulation
\usepackage{float} %used to properly place float-objects (figures)
\usepackage{titlesec, blindtext, color}
\usepackage{tabularx,ragged2e,booktabs,caption}
\usepackage{ulem} %used for strikeout
\usepackage{verbatim} %used for block-commenting
\usepackage[toc,page]{appendix}
\usepackage{url}
\usepackage{amsmath}
\usepackage{tikz}
\usetikzlibrary{shapes}

\definecolor{gray75}{gray}{0.75}
\definecolor{gray1}{gray}{0.97}
\definecolor{gray2}{gray}{0.90}
\definecolor{gray3}{gray}{0.80}
\definecolor{gray4}{gray}{0.63}

\lstset{
    basicstyle=\footnotesize\ttfamily, % Standardschrift
    %numbers=left,               % Ort der Zeilennummern
    numberstyle=\tiny,          % Stil der Zeilennummern
    %stepnumber=2,               % Abstand zwischen den Zeilennummern
    numbersep=5pt,              % Abstand der Nummern zum Text
    tabsize=2,                  % Groesse von Tabs
    extendedchars=\true,         %
    breaklines=true,            % Zeilen werden Umgebrochen
    keywordstyle=\color{red},
    frame=b,         
 %  keywordstyle=[1]\textbf,    % Stil der Keywords
 %  keywordstyle=[2]\textbf,    %
 %  keywordstyle=[3]\textbf,    %
 %  keywordstyle=[4]\textbf,   \sqrt{\sqrt{}} %
    % stringstyle=\color{white}\ttfamily, % Farbe der String
    showspaces=false,           % Leerzeichen anzeigen ?
    showtabs=false,       % Tabs anzeigen ?
    inputencoding=utf8,
    xleftmargin=17pt,
    framexleftmargin=17pt,
    framexrightmargin=5pt,
    framexbottommargin=4pt,
    %backgroundcolor=\color{lightgray},
    showstringspaces=false      % Leerzeichen in Strings anzeigen ? 
}


%\DeclareCaptionFont{blue}{\color{blue}} 

%\captionsetup[lstlisting]{singlelinecheck=false, labelfont={blue}, textfont={blue}}
\usepackage{caption}
\DeclareCaptionFont{white}{\color{white}}
\DeclareCaptionFormat{listing}{\colorbox[cmyk]{0.43, 0.35, 0.35,0.01}{\parbox{\textwidth}{\hspace{15pt}#1#2#3}}}
\captionsetup[lstlisting]{format=listing,labelfont=white,textfont=white, singlelinecheck=false, margin=0pt, font={bf,footnotesize}}


\newcommand{\hsp}{\hspace{20pt}}
\newcommand{\HRule}{\rule{\linewidth}{0.5mm}}



\newcommand{\img}[2]{ %images with borders
	\begin{figure}[H]
	\centering
	\fcolorbox{black}{black}{\includegraphics[width=14cm]{img/#1}}
	\caption[#2]{#2}
	\end{figure}
}

\newcommand{\imgC}[3]{ %images with borders and custom parameters
	\begin{figure}[H]
	\centering
	\fcolorbox{black}{black}{\includegraphics[#2]{img/#1}}
	\caption[#3]{#3}
	\end{figure}
}


\newcommand{\liste}[1]{\begin{enumerate}[label=\textbf{#1:\arabic*}]}
\newcommand{\inn}{\begin{enumerate}[label*=\textbf{.\arabic*}]}
\newcommand{\ut}{\end{enumerate}}


\newcommand{\refer}[5]{\bibitem {#1} #2. "#3" \textit{#4} (#5).}

%\newcommand{\labs}{\label{sec:utdypning}}
%\newcommand{\labl}[1]{\label{itm:#1}}
%\newcommand{\refs}{(se side~\pageref{sec:utdypning})}
%\newcommand{\refl}[1]{\ref{itm:#1} - }


%\parindent=5pt
%\baselineskip=0pt
%\parskip=0pt

\hypersetup{ %Brukes for � fjerne farger fra linker i content table
    colorlinks,
    citecolor=black,
    filecolor=black,
    linkcolor=black,
    urlcolor=black
}

\titleformat{\chapter}[hang]{\Huge\bfseries}{\thechapter\hsp\textcolor{gray75}{|}\hsp}{0pt}{\Huge\bfseries}




\begin{document}
\pagenumbering{Roman}
% Upper part of the page. The '~' is needed because \\
% only works if a paragraph has started.
\begin{titlepage}
\begin{center}
\includegraphics[width=0.15\textwidth]{img/ntnulogo.PNG}~\\[1cm]

\textsc{\LARGE NTNU}\\[1.5cm]

\textsc{\Large TDT4145 - Datamodellering og databasesystemer}\\[0.5cm]

% Title
\HRule \\[0.4cm]
{ \huge \bfseries Øving 3}\\[0.5cm]
{\large \textit{Mathias Ose og \O yvind Robertsen}}\\[0.2cm]
\HRule \\[1.5cm]



\vfill

% Bottom of the page
{\large \today}
\end{center}
\end{titlepage}

\newpage

\tableofcontents
\newpage

\pagenumbering{arabic}

\section{Oppgave 1}

\subsection{a)}

\textbf{Nøkkel:} Den enkelte eller den minimale mengden identifiserende attributtene ved en tabell. Må være unikt identifiserende for hver rad i tabellen. Kan være en enkelt attributt eller en mengde attributter.


\textbf{Supernøkkel:} En mengde attributter i en tabell som sammen fungerer som nøkkel for alle rader i tabellen. Kan også defineres som en mengde av attributter i en relasjon der alle andre attributter er funksjonelt avhengige av supernøkkelattributtene.


\textbf{Funksjonell avhengighet:} En funksjonell avhengighet $X \rightarrow Y$, der $X$ og $Y$ er mengder av attributter, betyr at verdiene av attributtene i $X$ bestemmer verdiene av $Y$.

\subsection{b)}

Tillukningen ($X^+$) til en mengde attributter $X$ er mengden av alle attributter i relasjonen R som er funksjonelt avhengig av $X$.

Algoritme for $X^+$, med hensyn på en mengde funksjonelle avhengigheter F:

\begin{lstlisting}[language=Python, caption=Tillukningsalgoritme, mathescape]
$X^+$ = $X$
$X^+_{gammel}$ = null
while not $X^+_{gammel} == X^+$:
$X^+_{gammel}$ = $X^+$
for functional dependency $Y \rightarrow Z$ in F:
if $Y$ in $X^+$:
$X^+$ = $X^+ \cup Z$
\end{lstlisting}

\subsection{c)}

\begin{gather*}
    a^+ = \{a, e\} \\
    ab^+ = \{a, b, c, d, e\} \\
    e^+ = \{e\}
\end{gather*}

\subsection{d)}

For at en mengde attributter skal være en supernøkkel, må tillukningen av mengden inneholde alle attributter i relasjonen. For en mengde attributter $X$ fra en relasjon $R$ avgjør vi om $X$ er en supernøkkel ved å finne $X^+$ og sjekke om $X^+$ inneholder alle attributter i $R$.

For at en supernøkkel også skal være en nøkkel, må den være \textit{minimal}. Dvs. at vi kan finne ut om en supernøkkel $SK$ også er en nøkkel ved å fjerne en og en attributt og sjekke om hver restmengde er en supernøkkel.

\subsection{e)}

For at projeksjonene $R_1$ og $R_2$ skal ha tapsløs-join-egenskapen må det gå an å joine $R_1$ og $R_2$ sammen og få $R$. Formelt har vi at $R_1$ og $R_2$ har tapsløs-join-egenskapen hvis og bare hvis

\begin{gather*}
(R_1 \cap R_2 \rightarrow R_1 - R_2 ) \in F^+ \vee (R_1 \cap R_2 \rightarrow R_2 - R_1) \in F^+
\end{gather*}

Det er naturlig å anta at det kun er mulig å ha tapsløs-join dersom $R_1 \cap R_2 \rightarrow R_1 \vee R_2$. Funksjonelle avhengigheter i R kan gi oss den formelle definisjonen over, der snittet mellom $R_1$ og $R_2$ ikke nødvendigvis er $R_1$ eller $R_2$. Dette på grunn av at funksjonelle avhengigheter kan gi oss muligheten til å manipulere snittet slik at det blir lik enten den ene eller den andre.

\subsection{f)}

Vi har relasjonen $R(a,b,c,d)$ med $F = \{b \rightarrow c\}$. Dette gir oss:

\begin{enumerate}
\item $R_1(a,b,c)$ og $R_2(b,c,d)$
    \begin{itemize}
        \item $R_1 \cap R_2 \rightarrow \{b,c\}$
        \item $R_1 - R_2 \rightarrow \{a\} \notin F^+$
        \item $R_2 - R_2 \rightarrow \{d\} \notin F^+$
        \item Ikke tapsløs join
    \end{itemize}
\item $R_1(a,b,d)$ og $R_2(b,c,d)$
    \begin{itemize}
        \item $R_1 \cap R_2 \rightarrow \{b,d\}$
        \item $R_1 - R_2 \rightarrow \{a\} \notin F^+$
        \item $R_2 - R-1 \rightarrow \{c\} \in F^+$
        \item Tapsløs join
    \end{itemize}
\item $R_1(a,b,d)$ og $R_2(b,c)$
    \begin{itemize}
    \item $R_1 \cap R_2 \rightarrow \{b\}$
    \item $R_1 - R_2 \rightarrow \{a,d\} \notin F^+$
    \item $R_2 - R_1 \rightarrow \{c\} \in F^+$
    \item Tapsløs join
    \end{itemize}
\end{enumerate}

\subsection{g)}

For at en relasjon $R$ skal være på tredje normalform, må følgende krav oppfylles:

\begin{itemize}
\item R er på andre normalform
\item Alle attributter som ikke er en del av en nøkkel, er direkte avhengig av hele nøkkelen. Dvs. at transitiv avhengighet ikke er lov
\end{itemize}

\subsection{h)}

Vi har en relasjon $R(A,B,C,D)$ med funksjonelle avhengigheter $F = \{A \rightarrow B, C \rightarrow D\}$ som ikke er på 3NF. For å dekomponere til relasjoner som oppfyller kravene for 3NF, gjør vi følgende:

\begin{enumerate}
\item Finner et såkalt canonical cover $F_C$ for $F$. (I dette tilfellet er $F_C = F$)
\item For hver FD $X \rightarrow A \in F_C$, lager vi en relasjon $XA$ med $X$ som nøkkel.
    \begin{itemize}
    \item For oss blir dette $R_1(A,B)$ med $A$ som nøkkel og $R_2(C,D)$ med $C$ som nøkkel.
    \end{itemize}
\item De resterende attributtene og nøkkel samler vi i en siste relasjon. I vårt tilfelle inneholder denne relasjonen kun nøkkelen fra den opprinnelige relasjonen, nemlig $R_3(A,C)$.
\end{enumerate}

Ser at denne fremgangsmåten er fulgt for å komme fra til dekomponeringen i oppgaven. Ettersom denne fremgangsmåten bevarer funksjonelle avhengigheter og garanterer lossless-join, er dette et godt valg.

\section{Oppgave 2}

\newpage
\section{Oppgave 3}
~\\
Setter inn verdien $12$: $12 \mod 5 = 2$\\
\begin{tabular}{|l|l|l|}
    \hline
    Adresse & Post 1 & Post 2 \\
    0       & ~      & ~      \\
    1       & ~      & ~      \\
    2       & 12     & ~      \\
    3       & ~      & ~      \\
    4       & ~      & ~      \\ \hline
\end{tabular}

~\\
Setter inn verdien $9$: $9 \mod 5 = 4$\\
\begin{tabular}{|l|l|l|}
    \hline
    Adresse & Post 1 & Post 2 \\
    0       & ~      & ~      \\
    1       & ~      & ~      \\
    2       & 12     & ~      \\
    3       & ~      & ~      \\
    4       & 9      & ~      \\ \hline
\end{tabular}

~\\
Setter inn verdien $3$: $3 \mod 5 = 3$\\
\begin{tabular}{|l|l|l|}
    \hline
    Adresse & Post 1 & Post 2 \\
    0       & ~      & ~      \\
    1       & ~      & ~      \\
    2       & 12     & ~      \\
    3       & 3      & ~      \\
    4       & 9      & ~      \\ \hline
\end{tabular}

~\\
Setter inn verdien $18$: $18 \mod 5 = 3$\\
Her må altså post 2 tas i bruk.\\
\begin{tabular}{|l|l|l|}
    \hline
    Adresse & Post 1 & Post 2 \\
    0       & ~      & ~      \\
    1       & ~      & ~      \\
    2       & 12     & ~      \\
    3       & 3      & 18      \\
    4       & 9      & ~      \\ \hline
\end{tabular}

~\\
Setter inn verdien $22$: $22 \mod 5 = 2$\\
\begin{tabular}{|l|l|l|}
    \hline
    Adresse & Post 1 & Post 2 \\
    0       & ~      & ~      \\
    1       & ~      & ~      \\
    2       & 12     & 22     \\
    3       & 3      & 18     \\
    4       & 9      & ~      \\ \hline
\end{tabular}

~\\
Setter inn verdien $1$: $1 \mod 5 = 1$\\
\begin{tabular}{|l|l|l|}
    \hline
    Adresse & Post 1 & Post 2 \\
    0       & ~      & ~      \\
    1       & 1      & ~      \\
    2       & 12     & 22     \\
    3       & 3      & 18     \\
    4       & 9      & ~      \\ \hline
\end{tabular}

~\\
Setter inn verdien $37$: $37 \mod 5 = 2$\\
Det finnes ingen ledig post for verdien, så vi oppretter en overflyts-blokk med to nye poster.\\
\begin{tabular}{|l|l|l|l|l|}
    \hline
    Adresse & Post 1 & Post 2 & Of1 p1 & Of1 p2 \\
    0       & ~      & ~      &\cellcolor{gray}&\cellcolor{gray}\\
    1       & 1      & ~      &\cellcolor{gray}&\cellcolor{gray}\\
    2       & 12     & 22     & 37     & ~      \\
    3       & 3      & 18     &\cellcolor{gray}&\cellcolor{gray}\\
    4       & 9      & ~      &\cellcolor{gray}&\cellcolor{gray}\\ \hline
\end{tabular}

~\\
Setter inn verdien $11$: $11 \mod 5 = 1$\\
\begin{tabular}{|l|l|l|l|l|}
    \hline
    Adresse & Post 1 & Post 2 & Of1 p1 & Of1 p2 \\
    0       & ~      & ~      &\cellcolor{gray}&\cellcolor{gray}\\
    1       & 1      & 11     &\cellcolor{gray}&\cellcolor{gray}\\
    2       & 12     & 22     & 37     & ~      \\
    3       & 3      & 18     &\cellcolor{gray}&\cellcolor{gray}\\
    4       & 9      & ~      &\cellcolor{gray}&\cellcolor{gray}\\ \hline
\end{tabular}

~\\
Setter inn verdien $87$: $87 \mod 5 = 2$\\
\begin{tabular}{|l|l|l|l|l|}
    \hline
    Adresse & Post 1 & Post 2 & Of1 p1 & Of1 p2 \\
    0       & ~      & ~      &\cellcolor{gray}&\cellcolor{gray}\\
    1       & 1      & 11     &\cellcolor{gray}&\cellcolor{gray}\\
    2       & 12     & 22     & 37     & 87     \\
    3       & 3      & 18     &\cellcolor{gray}&\cellcolor{gray}\\
    4       & 9      & ~      &\cellcolor{gray}&\cellcolor{gray}\\ \hline
\end{tabular}

~\\
Setter inn verdien $2$: $2 \mod 5 = 2$\\
Det finnes ingen ledig post for verdien, så vi må opprette en ny overflyts-blokk.\\
\begin{tabular}{|l|l|l|l|l|l|l|}
    \hline
    Adresse & Post 1 & Post 2 & Of1 p1 & Of1 p2 & Of2 p1 & Of2 p2 \\
    0       & ~      & ~      &\cellcolor{gray}&\cellcolor{gray}&\cellcolor{gray}&\cellcolor{gray}\\
    1       & 1      & 11     &\cellcolor{gray}&\cellcolor{gray}&\cellcolor{gray}&\cellcolor{gray}\\
    2       & 12     & 22     & 37     & 87     & 2      & ~      \\
    3       & 3      & 18     &\cellcolor{gray}&\cellcolor{gray}&\cellcolor{gray}&\cellcolor{gray}\\
    4       & 9      & ~      &\cellcolor{gray}&\cellcolor{gray}&\cellcolor{gray}&\cellcolor{gray}\\ \hline
\end{tabular}

~\\
Setter inn verdien $1$: $1 \mod 5 = 1$\\
Det finnes ingen ledig post for verdien, så vi må opprette en ny overflyts-blokk.\\
\begin{tabular}{|l|l|l|l|l|l|l|}
    \hline
    Adresse & Post 1 & Post 2 & Of1 p1 & Of1 p2 & Of2 p1 & Of2 p2 \\
    0       & ~      & ~      &\cellcolor{gray}&\cellcolor{gray}&\cellcolor{gray}&\cellcolor{gray}\\
    1       & 1      & 11     & 1      & ~      &\cellcolor{gray}&\cellcolor{gray}\\
    2       & 12     & 22     & 37     & 87     & 2      & ~      \\
    3       & 3      & 18     &\cellcolor{gray}&\cellcolor{gray}&\cellcolor{gray}&\cellcolor{gray}\\
    4       & 9      & ~      &\cellcolor{gray}&\cellcolor{gray}&\cellcolor{gray}&\cellcolor{gray}\\ \hline
\end{tabular}


\end{document}
